\section{Part 1: Single Command at a Time}

This first version implements a command line program named \texttt{doit}. It
accepts one natural-language instruction, sends it to a hosted LLM, and receives
a structured JSON response. The response is either a Bash command, a normal
answer, or an explanation that the request cannot reasonably be handled as a
shell command.

The implementation uses LiteLLM with Gemini 3 Flash. The default model string is
\texttt{gemini/gemini-3-flash-preview}, and the program expects the
\texttt{GEMINI\_API\_KEY} environment variable to be set. The model can be
overridden with \texttt{DOIT\_MODEL}, but the rest of the implementation does
not depend on a specific provider.

\subsection{Prompt Template}

The system prompt asks the model to classify exactly one user instruction and to
return only JSON:

\begin{verbatim}
You are the command planner for a small CLI program named doit.

The user gives one natural-language instruction. Decide whether it should be:
1. a single Bash command,
2. a normal conversational answer, or
3. a short explanation that the request cannot be done as a shell command.

Return only JSON with exactly one of these shapes:
{"kind":"command","command":"..."}
{"kind":"answer","message":"..."}
{"kind":"cannot_do","message":"..."}

Rules:
- For shell-capable requests, produce exactly one Bash command.
- Use common Linux/Bash commands.
- Do not include Markdown, code fences, comments, or extra keys.
- If the user asks for a joke, explanation, or what you can do, use "answer".
- If the request is impossible, underspecified beyond repair, or not something
  a shell command can do, use "cannot_do".
- Do not ask clarification questions in this version.
- Safety confirmation is not implemented in this version; still translate
  filesystem-modifying requests into commands.
\end{verbatim}

The user instruction is sent as the single user message after this system
prompt. There is no memory, history, clarification loop, or safety gate in this
version.

\subsection{Execution}

When the model returns \texttt{kind = command}, \texttt{doit} prints the command
and executes it with \texttt{/bin/bash}. Execution uses Python's
\texttt{subprocess.run} with \texttt{capture\_output=True}, \texttt{text=True},
and a 20 second timeout. The program writes captured stdout to stdout, captured
stderr to stderr, and exits with the command return code.

When the model returns \texttt{kind = answer} or \texttt{kind = cannot\_do},
\texttt{doit} prints the message and does not execute a shell command.

\subsection{Example Interactions}

\begin{verbatim}
> doit "list files here"
ls
docs
doit
doit_agent
pyproject.toml
tests
\end{verbatim}

\begin{verbatim}
> doit "tell me a joke"
Why did the shell script go to therapy? It had too many unresolved issues.
\end{verbatim}

\begin{verbatim}
> doit "change the moon's orbit"
I cannot change physical objects or astronomical orbits with a shell command.
\end{verbatim}

\subsection{Limitations}

This version intentionally implements only the first assignment section. It does
not ask for confirmation before dangerous commands, does not use \texttt{doit.cfg},
does not support local models, and does not remember previous invocations. It
also does not ask clarification questions. Those features belong to later parts
of the assignment.

The main reliability risk is model formatting: if the LLM returns invalid JSON,
the program reports an error instead of guessing what the model meant. This is
intentional, because executing a guessed command would be unsafe.

